\input{preamble/preamble.tex}

\geometry{margin=0.85in}
\AtBeginDocument{\small}

\newcommand{\ExamNameField}{\noindent\textbf{Name:}\ \rule{0.7\linewidth}{0.4pt}\par\medskip}

\begin{document}
	\ExamTitleBlock{11th grade}{Term 2 -- Lesson 3 MATERIAL: t-Confidence Intervals for Means (Worked Solutions)}{}
	
	\ExamSection{Evaluated criteria}
	\begin{ExamCriteria}
		\item C8: Calculates confidence intervals applying the t-student distribution in context situations.
		\item C9: Estimates the confidence interval for the mean when the variance is unknown in context situations.
		\item C10: Concludes about a statistical parameter in situations in finance and economy.
	\end{ExamCriteria}

	\ExamSection{C8: t-student confidence intervals in context situations}
	\begin{ExamProblems}
		\item
		\subsection*{Problem description}
		A logistics center wants to estimate the mean loading time (minutes) for trucks in a new dispatch lane.
		A random sample of $n=12$ trucks gives $\bar{x}=84.60$ minutes and sample standard deviation $s=6.80$ minutes.
		The population variance is unknown, and the sample size is small, so a t-interval is required.
		Construct and interpret a 95\% confidence interval for the population mean loading time.
		
		\subsection*{C8}
		Step 1.
		Identify the parameter and sample statistics.
		\[
		\mu=\text{population mean loading time},\quad n=12,\quad \bar{x}=84.60,\quad s=6.80
		\]
		Step 2.
		State degrees of freedom.
		\[
		df=n-1=12-1=11
		\]
		Step 3.
		Identify the t critical value for 95\% confidence.
		\[
		t^*=t_{0.025,11}\approx 2.201
		\]
		Step 4.
		Compute the standard error.
		\[
		SE=\frac{s}{\sqrt{n}}=\frac{6.80}{\sqrt{12}}\approx 1.96
		\]
		Step 5.
		Compute the margin of error.
		\[
		E=t^*\cdot SE=2.201(1.96)\approx 4.31
		\]
		Step 6.
		Construct the confidence interval.
		\[
		\mu\in\bar{x}\pm E\Rightarrow 84.60\pm 4.31=[80.29,\,88.91]
		\]
		Step 7.
		Interpretation.
		With 95\% confidence, the mean loading time in the new lane is between 80.29 and 88.91 minutes.
		
		\newpage
		\item
		\subsection*{Problem description}
		An agricultural exporter tracks the average customs processing time (hours) for a new shipping route.
		A random sample of $n=18$ shipments gives $\bar{x}=42.30$ hours and $s=5.10$ hours.
		Because the population variance is unknown and the sample is moderate, use a t-interval.
		Construct and interpret a 90\% confidence interval for the population mean processing time.
		
		\subsection*{C8}
		Step 1.
		Identify the parameter and sample statistics.
		\[
		\mu=\text{population mean processing time},\quad n=18,\quad \bar{x}=42.30,\quad s=5.10
		\]
		Step 2.
		State degrees of freedom.
		\[
		df=18-1=17
		\]
		Step 3.
		Identify the t critical value for 90\% confidence.
		\[
		t^*=t_{0.05,17}\approx 1.740
		\]
		Step 4.
		Compute the standard error.
		\[
		SE=\frac{5.10}{\sqrt{18}}\approx 1.20
		\]
		Step 5.
		Compute the margin of error.
		\[
		E=1.740(1.20)\approx 2.09
		\]
		Step 6.
		Construct the confidence interval.
		\[
		\mu\in 42.30\pm 2.09=[40.21,\,44.39]
		\]
		Step 7.
		Interpretation.
		With 90\% confidence, the true mean customs processing time for this route lies between 40.21 and 44.39 hours.
		
		\newpage
		\item
		\subsection*{Problem description}
		A school cafeteria estimates the mean daily demand (kg) for a new meal plan ingredient.
		A random sample of $n=25$ school days gives $\bar{x}=73.50$ kg and $s=8.40$ kg.
		The population variance is unknown.
		Construct and interpret a 99\% confidence interval for the population mean daily demand.
		
		\subsection*{C8}
		Step 1.
		Identify the parameter and sample statistics.
		\[
		\mu=\text{population mean daily demand},\quad n=25,\quad \bar{x}=73.50,\quad s=8.40
		\]
		Step 2.
		State degrees of freedom.
		\[
		df=25-1=24
		\]
		Step 3.
		Identify the t critical value for 99\% confidence.
		\[
		t^*=t_{0.005,24}\approx 2.797
		\]
		Step 4.
		Compute the standard error.
		\[
		SE=\frac{8.40}{\sqrt{25}}=1.68
		\]
		Step 5.
		Compute the margin of error.
		\[
		E=2.797(1.68)\approx 4.70
		\]
		Step 6.
		Construct the confidence interval.
		\[
		\mu\in 73.50\pm 4.70=[68.80,\,78.20]
		\]
		Step 7.
		Interpretation.
		With 99\% confidence, the population mean daily demand is between 68.80 kg and 78.20 kg.
	\end{ExamProblems}

	\ExamSection{C9: Confidence interval for the mean when variance is unknown}
	\begin{ExamProblems}
		\item
		\subsection*{Problem description}
		A customer-service unit studies the average response time (minutes) for premium support tickets.
		The population variance is unknown.
		A random sample of 10 tickets gives the following response times:
		\[
		14,\ 16,\ 15,\ 13,\ 17,\ 18,\ 16,\ 14,\ 15,\ 17
		\]
		Construct and interpret a 95\% confidence interval for the population mean response time.
		
		\subsection*{C9}
		Step 1.
		Compute the sample mean.
		\[
		\bar{x}=\frac{14+16+15+13+17+18+16+14+15+17}{10}=\frac{155}{10}=15.50
		\]
		Step 2.
		Compute the sample standard deviation.
		\[
		s^2=\frac{\sum (x_i-\bar{x})^2}{n-1}=\frac{22.50}{9}=2.50,
		\qquad s=\sqrt{2.50}\approx 1.58
		\]
		Step 3.
		State degrees of freedom and the critical value.
		\[
		df=10-1=9,\qquad t^*=t_{0.025,9}\approx 2.262
		\]
		Step 4.
		Compute the standard error.
		\[
		SE=\frac{s}{\sqrt{n}}=\frac{1.58}{\sqrt{10}}\approx 0.50
		\]
		Step 5.
		Compute the margin of error.
		\[
		E=t^*\cdot SE=2.262(0.50)\approx 1.13
		\]
		Step 6.
		Construct the confidence interval.
		\[
		\mu\in 15.50\pm 1.13=[14.37,\,16.63]
		\]
		Step 7.
		Interpretation.
		With 95\% confidence, the mean response time for premium tickets lies between 14.37 and 16.63 minutes.
		
		\newpage
		\item
		\subsection*{Problem description}
		A local retail store estimates mean daily sales (thousand USD) for a new product line.
		The population variance is unknown.
		A random sample of 8 days gives:
		\[
		52,\ 49,\ 55,\ 47,\ 50,\ 54,\ 51,\ 48
		\]
		Construct and interpret a 90\% confidence interval for the population mean daily sales.
		
		\subsection*{C9}
		Step 1.
		Compute the sample mean.
		\[
		\bar{x}=\frac{52+49+55+47+50+54+51+48}{8}=\frac{406}{8}=50.75
		\]
		Step 2.
		Compute the sample standard deviation.
		\[
		s^2=\frac{\sum (x_i-\bar{x})^2}{n-1}=\frac{55.50}{7}\approx 7.93,
		\qquad s\approx \sqrt{7.93}=2.82
		\]
		Step 3.
		State degrees of freedom and the critical value.
		\[
		df=8-1=7,\qquad t^*=t_{0.05,7}\approx 1.895
		\]
		Step 4.
		Compute the standard error.
		\[
		SE=\frac{2.82}{\sqrt{8}}\approx 1.00
		\]
		Step 5.
		Compute the margin of error.
		\[
		E=1.895(1.00)\approx 1.90
		\]
		Step 6.
		Construct the confidence interval.
		\[
		\mu\in 50.75\pm 1.90=[48.85,\,52.65]
		\]
		Step 7.
		Interpretation.
		With 90\% confidence, the population mean daily sales are between 48.85 and 52.65 thousand USD.
		
		\newpage
		\item
		\subsection*{Problem description}
		An energy analyst estimates mean monthly electricity cost (USD) for small stores in one district.
		The population variance is unknown.
		A random sample of 12 stores gives:
		\[
		31,\ 29,\ 34,\ 33,\ 30,\ 32,\ 35,\ 28,\ 31,\ 33,\ 30,\ 34
		\]
		Construct and interpret a 95\% confidence interval for the population mean monthly cost.
		
		\subsection*{C9}
		Step 1.
		Compute the sample mean.
		\[
		\bar{x}=\frac{380}{12}\approx 31.67
		\]
		Step 2.
		Compute the sample standard deviation.
		\[
		s^2=\frac{52.67}{11}\approx 4.79,
		\qquad s\approx\sqrt{4.79}=2.19
		\]
		Step 3.
		State degrees of freedom and the critical value.
		\[
		df=12-1=11,\qquad t^*=t_{0.025,11}\approx 2.201
		\]
		Step 4.
		Compute the standard error.
		\[
		SE=\frac{2.19}{\sqrt{12}}\approx 0.63
		\]
		Step 5.
		Compute the margin of error.
		\[
		E=2.201(0.63)\approx 1.39
		\]
		Step 6.
		Construct the confidence interval.
		\[
		\mu\in 31.67\pm 1.39=[30.28,\,33.06]
		\]
		Step 7.
		Interpretation.
		With 95\% confidence, the district mean monthly electricity cost for small stores is between 30.28 and 33.06 USD.
	\end{ExamProblems}

	\ExamSection{C10: Conclusions about a mean parameter in finance and economy}
	\begin{ExamProblems}
		\item
		\subsection*{Problem description}
		A microfinance institution evaluates the mean approval time (days) for small-business loans.
		Managers claim the true mean approval time is 4.0 days.
		A random sample gives $n=15$, $\bar{x}=4.80$ days, and $s=1.20$ days.
		The population variance is unknown.
		Construct a 95\% t-confidence interval for the population mean approval time and conclude whether the 4.0-day claim is plausible.
		
		\subsection*{C10}
		Step 1.
		Identify the parameter and sample statistics.
		\[
		\mu=\text{population mean approval time},\quad n=15,\quad \bar{x}=4.80,\quad s=1.20
		\]
		Step 2.
		State degrees of freedom and critical value.
		\[
		df=15-1=14,\qquad t^*=t_{0.025,14}\approx 2.145
		\]
		Step 3.
		Compute the standard error.
		\[
		SE=\frac{1.20}{\sqrt{15}}\approx 0.31
		\]
		Step 4.
		Compute the margin of error.
		\[
		E=2.145(0.31)\approx 0.66
		\]
		Step 5.
		Construct the confidence interval.
		\[
		\mu\in 4.80\pm 0.66=[4.14,\,5.46]
		\]
		Step 6.
		Interpretation in context.
		With 95\% confidence, the institution's mean approval time is between 4.14 and 5.46 days.
		Step 7.
		Conclusion about the parameter.
		Because 4.0 days is outside the interval, the managerial claim is not supported by this sample evidence.
		
		\newpage
		\item
		\subsection*{Problem description}
		An investment team estimates the mean daily return (\%) of a balanced fund.
		A promotional report states that the true mean return is 0.30\% per day.
		From a random sample of $n=20$ trading days, $\bar{x}=0.42\%$ and $s=0.18\%$.
		The population variance is unknown.
		Construct a 90\% t-confidence interval and conclude whether the reported value is plausible.
		
		\subsection*{C10}
		Step 1.
		Identify the parameter and sample statistics.
		\[
		\mu=\text{population mean daily return},\quad n=20,\quad \bar{x}=0.42\%,\quad s=0.18\%
		\]
		Step 2.
		State degrees of freedom and critical value.
		\[
		df=20-1=19,\qquad t^*=t_{0.05,19}\approx 1.729
		\]
		Step 3.
		Compute the standard error.
		\[
		SE=\frac{0.18\%}{\sqrt{20}}\approx 0.04\%
		\]
		Step 4.
		Compute the margin of error.
		\[
		E=1.729(0.04\%)\approx 0.07\%
		\]
		Step 5.
		Construct the confidence interval.
		\[
		\mu\in 0.42\%\pm 0.07\%=[0.35\%,\,0.49\%]
		\]
		Step 6.
		Interpretation in context.
		With 90\% confidence, the mean daily return is between 0.35\% and 0.49\%.
		Step 7.
		Conclusion about the parameter.
		Since 0.30\% is outside the interval, the reported mean appears too low for this sample period.
		
		\newpage
		\item
		\subsection*{Problem description}
		A digital payments company estimates the mean transaction fee (USD) charged to merchants.
		The pricing team states that the long-run mean fee is 2.20 USD and should remain below 2.40 USD.
		A random sample of $n=9$ transactions gives $\bar{x}=2.35$ USD and $s=0.41$ USD.
		The population variance is unknown.
		Construct a 99\% t-confidence interval for the population mean fee and conclude about the pricing statement.
		
		\subsection*{C10}
		Step 1.
		Identify the parameter and sample statistics.
		\[
		\mu=\text{population mean fee},\quad n=9,\quad \bar{x}=2.35,\quad s=0.41
		\]
		Step 2.
		State degrees of freedom and critical value.
		\[
		df=9-1=8,\qquad t^*=t_{0.005,8}\approx 3.355
		\]
		Step 3.
		Compute the standard error.
		\[
		SE=\frac{0.41}{\sqrt{9}}=0.14
		\]
		Step 4.
		Compute the margin of error.
		\[
		E=3.355(0.14)\approx 0.46
		\]
		Step 5.
		Construct the confidence interval.
		\[
		\mu\in 2.35\pm 0.46=[1.89,\,2.81]
		\]
		Step 6.
		Interpretation in context.
		With 99\% confidence, the population mean transaction fee lies between 1.89 and 2.81 USD.
		Step 7.
		Conclusion about the parameter.
		The value 2.20 USD is plausible because it lies inside the interval, but the interval also includes values above 2.40 USD, so the ``below 2.40 USD'' condition is not guaranteed at 99\% confidence with this sample.
	\end{ExamProblems}
\end{document}
